\subsection{可分代数元素}

定义3.——设E是K的一个扩域。E中在K上代数的元素x称为在K上可分，如果代数扩张 \(\mathrm{K}(x)\) 是K的可分扩张。

命题5.——设E是K的一个扩域，x是E中在K上代数的元素，f是x在K上的极小多项式。则以下条件等价：

\begin{enumerate}
\def\labelenumi{\alph{enumi})}
\item
  \(x\) （在 \(\mathbf{K}\) 上）是可分的；
\item
  多项式 \(f\) 是可分的；
\item
  \(x\) 是 \(f\) 的一个单根。
\end{enumerate}

\(a)\) 与 \(b)\) 的等价性由命题3得出，\(b)\) 与 \(c)\) 的等价性由命题3和4得出（参见V，第37和38页）。

推论1. --- 若 E 的元素 x 是 K{[}X{]} 中多项式 g 的单根，则它在 K 上是可分的。

由于极小多项式 \(f\) （\(x\) 在 \(K\) 上的）整除 \(g\) （在 \(K[X]\) 中）（V，第16页，定理1），所以 \(x\) 是 \(f\) 的一个单根。

推论2. —— 若E的元素x在K上是代数的且可分的，则它在E所包含的K的每个扩张K'上也是代数的且可分的。

设 \(f\) 为 \(x\) 在 \(K\) 上的极小多项式，则由命题5，\(x\) 是 \(f\) 的一个单根，且由于 \(f\) 属于 \(K'[X]\) ，元素 \(x\) （\(E\) 的）在 \(K'\) 上是可分的（由推论1）。

推论3. --- 假设 K 的特征为 \(p \neq 0\) 。要使 E 的元素 \(x\) 属于 K，必要且充分的是它在 K 上既是可分代数的又是 \(p\) -根式的。

所述条件显然是必要的。反之，假设 \(x\) 在 \(\mathbf{K}\) 上是可分代数的，并且是 \(p\) -根式的，高度为 \(e\) （在 \(\mathbf{K}\) 上）。由于 \(x\) 在 \(\mathbf{K}\) 上可分，极小多项式 \(f\) （\(x\) 在 \(\mathbf{K}\) 上的）不属于 \(\mathbf{K}[X^p]\) （命题4和5）；由于 \(x\) 是 \(p\) -根式的，高度为 \(e\) （在 \(\mathbf{K}\) 上），我们有 \(f(\mathbf{X}) = \mathbf{X}^{p^e} - x^{p^e}(\mathbf{V},\text{p. 24, Prop. 1})\) ；我们得出结论 \(e = 0\) ，所以 \(x \in \mathbf{K}\) 。

命题6. --- 设E是K的一个扩域。

\begin{enumerate}
\def\labelenumi{\alph{enumi})}
\item
  如果 \(\mathbf{E}\) 在 \(\mathbf{K}\) 上是代数的且可分的，则 \(\mathbf{E}\) 的每个元素都在 \(\mathbf{K}\) 上是代数的且可分的。
\item
  反之，设 \(\mathbf{A}\) 是 \(\mathbf{E}\) 的一个元素集合，这些元素在 \(\mathbf{K}\) 上是代数的且可分的，并且满足 \(\mathbf{E} = \mathbf{K}(\mathbf{A})\) ；则 \(\mathbf{E}\) 在 \(\mathbf{K}\) 上是代数的且可分的。
\end{enumerate}

如果 \(\mathbf{E}\) 在 \(\mathbf{K}\) 上是代数的且可分的，则扩张 \(\mathbf{K}(x)\) （\(\mathbf{K}\) 的）对每个 \(x \in \mathbf{E}\) 也如此，从而得 \(a\) ).

在假设 \(b\) ) 下，扩张 \(\mathbf{E}\) 在 \(\mathbf{K}\) 上是代数的（V，第18页，推论1）。

设 \(\mathbf{F}\) 是 \(\mathbf{E}\) 的一个在 \(\mathbf{K}\) 上有限次的子扩张。根据 \(\mathbf{V}\) 第～11页推论，存在元素 \(x_{1},\ldots ,x_{m}\) （属于 \(\mathbf{A}\)），使得 \(\mathbf{F}\subset \mathbf{K}(x_1,\dots,x_m)\) ，并且我们有

\[
\mathrm{K} \left(x _ {1}, \dots , x _ {m}\right) = \mathrm{K} \left[ x _ {1}, \dots , x _ {m} \right] \quad (\mathrm{V}, \text { p. } 1 8, \text { Cor. } 1).
\]

根据对 A 的假设，代数 \(K[x_{1}], \ldots, K[x_{m}]\) 在 K 上是平展的；因此 \(K[x_{1}] \otimes \cdots \otimes K[x_{m}]\) 同样如此（V，第32页，推论1）。现在 F 同构于 \(K[x_{1}] \otimes \cdots \otimes K[x_{m}]\) 的某个商代数的一个子代数，因此是平展的（V，第30页，命题3）。

推论。—— 要使代数扩张 E 在 K 上可分，必要且充分的是 E 的每个元素都是它在 K 上的极小多项式的单根。只需应用命题5和命题6即可。

